% =====================================================================
%  TRABAJOS FUTUROS
%  2026-10-15: al día con lo hecho (la unidad inercial y los dos vehículos
%  en la misma red ya no son trabajo futuro), sin códigos del repositorio,
%  y con dos líneas nuevas apoyadas en la literatura del marco teórico.
%  2026-10-19: cerrada con lo que dejó la campaña de los dos vehículos
%  (medida externa de la llegada y estabilidad en sesiones largas).
% =====================================================================
\chapter*{Trabajos futuros}
\addcontentsline{toc}{chapter}{TRABAJOS FUTUROS}

Las líneas que siguen salen de límites concretos que el trabajo encontró y
documentó, y de las suposiciones que el diseño hizo para mantenerse dentro del
alcance.

\paragraph{Aproximación fina del vehículo.} Es la limitación que condiciona la
precisión de llegada, y está localizada y medida
(sección~\ref{sec:banda-muerta}). El vehículo no se mueve por debajo de
0,40~m/s, porque por debajo de ese valor la cadena de tracción no entrega
potencia suficiente. Hay dos vías para resolverla. Una es de
hardware: calibrar el variador o cambiar la relación de transmisión, para darle al
vehículo un régimen lento del que hoy carece. La otra es de control: cerrar el
lazo de velocidad con una medida de avance, de modo que el vehículo pueda
regularla cerca de la meta.

\paragraph{Una fuente de avance independiente del láser.} La unidad inercial
resolvió el rumbo, pero la posición sigue saliendo de comparar barridos del
láser, y eso falla por construcción en pasillos uniformes
(sección~\ref{sec:teo-localizacion}). El proyecto lo resolvió eligiendo un entorno
con estructura, y esa solución limita dónde puede operar el sistema.
Codificadores en las ruedas, u odometría visual con las cámaras que los vehículos
ya llevan, darían información de avance que no depende de la geometría del sitio.

\paragraph{Una medida externa de la llegada.} En la campaña con los dos
vehículos la llegada se verificó con la localización del propio vehículo
(sección~\ref{sec:lim-llegada}). Repetirla con una medida independiente daría la
distancia real a cada destino. Puede ser el flexómetro sobre marcas en el piso,
como pedía el protocolo, o un sistema de seguimiento externo, que además
registraría la trayectoria completa sin depender del vehículo.

\paragraph{Estabilidad de la navegación en sesiones largas.} Tras horas
encendidos, la navegación de los dos vehículos dejó de responder y hubo que
reiniciarla (sección~\ref{sec:res-completo}). La causa no se estableció. Antes de
dejar el sistema en servicio hace falta reproducir la falla en una sesión larga
controlada. Con ella se puede decidir si basta con supervisar la cadena de
navegación y reiniciarla sola, o si hay que bajar la carga del procesador del
vehículo.

\paragraph{El ascensor como segunda vía de transición.} Hoy el cambio de piso se
hace siempre por la escalera, y el ascensor es un destino más. Admitir los dos
exige que el planificador elija entre dos puntos de transferencia por nivel,
presumiblemente el más cercano al usuario o el que este prefiera. El planificador
ya detecta ese caso y falla con un mensaje explícito cuando encuentra dos, de
modo que el punto de extensión está identificado.

\paragraph{Esquinas en pasillos estrechos.} Todos los recorridos medidos fueron
sobre tramos rectos. La media vuelta en varios tiempos resolvió la inversión de
sentido en pasillos de 2,5~m. Falta probar que el vehículo resuelva un giro de
90$^\circ$ en una esquina estrecha, que es lo que exige el guiado en edificios con
más de un pasillo por piso.

\paragraph{Más de un vehículo por nivel y misiones simultáneas.} El sistema
atiende una misión a la vez, y cada nivel tiene un solo vehículo. Con esas dos
condiciones la asignación es trivial (sección~\ref{sec:teo-asignacion}). En cuanto
haya dos o más vehículos elegibles por tarea, o varias solicitudes a la vez, la
asignación deja de ser trivial, y los métodos de subasta, como
MURDOCH \autocite{gerkey2002sold} o CBBA \autocite{choi2009consensus}, empiezan a
decidir algo. Es la extensión natural del trabajo para edificios con más
afluencia.

\paragraph{Ampliar la muestra de las campañas.} Con treinta misiones en
simulación el intervalo de confianza mide 24 puntos porcentuales, y con nueve en
los vehículos la cota inferior es del 70~\%. Ese ancho es de diseño y solo se
reduce con más repeticiones. Una campaña mayor permitiría además comparar las
misiones de un mismo piso con las de dos pisos, comparación que hoy no se sostiene
porque sus intervalos se solapan.

\paragraph{Evaluación con usuarios.} Las cuatro métricas del proyecto miden el
sistema, y no la experiencia de quien lo usa. Un estudio con personas, como el
que hicieron \textcite{guerreiro2019cabot} con su robot guía, permitiría medirla.
Interesa saber si la interfaz se entiende sin explicación, si la velocidad del
vehículo resulta cómoda para caminar a su lado y cuánto tolera el usuario la
espera en la escalera. Esa
evaluación es también el paso previo a ofrecer el sistema a personas con
discapacidad visual, que es la aplicación con más impacto que sugiere la
literatura revisada.
